% =====================================================================
% Supplementary Information for paper06_assessment_separation_v77
% =====================================================================
\documentclass[preprint,11pt,authoryear]{elsarticle}
\usepackage{amsmath,amssymb,amsthm}
\usepackage{booktabs}
\usepackage{array}
\usepackage{tabularx}
\usepackage{graphicx}
\graphicspath{{./}{figs/}{paper06_assessment_separation_v77_src/figs/}}
\usepackage{enumitem}
\usepackage[colorlinks=true,urlcolor=blue,linkcolor=blue,citecolor=blue]{hyperref}
\usepackage{xr}
\externaldocument[main-]{paper06_assessment_separation_v77}
\emergencystretch=3em
\setlength{\tolerance}{2000}

\renewcommand{\thesection}{S\arabic{section}}
\renewcommand{\thetable}{S\arabic{table}}
\renewcommand{\thefigure}{S\arabic{figure}}
\renewcommand{\theequation}{S\arabic{equation}}
\newtheorem{theorem}{Theorem}
\renewcommand{\thetheorem}{S\arabic{theorem}}
\newtheorem{lemma}{Lemma}
\renewcommand{\thelemma}{S\arabic{lemma}}
\newtheorem{proposition}{Proposition}
\renewcommand{\theproposition}{S\arabic{proposition}}
\theoremstyle{remark}
\newtheorem{remark}{Remark}
\renewcommand{\theremark}{S\arabic{remark}}

\newcommand{\Etyp}{E_{\mathrm{typ}}}
\newcommand{\Eend}{E_{\mathrm{end}}}
\newcommand{\Ephys}{E_{\mathrm{tube,phys}}}
\newcommand{\Vtyp}{\mathcal{V}_{\mathrm{typ}}}
\newcommand{\Vweak}{\mathcal{V}_{\mathrm{weak}}}
\newcommand{\Vphys}{\mathcal{V}_{\mathrm{phys}}}
\newcommand{\Vend}{\mathcal{V}_{\mathrm{end}}}
\newcommand{\Wp}{W_+}
\newcommand{\FP}{\mathrm{FP}_{\mathrm{agg}}}
\newcommand{\Tube}{\mathrm{Tube}}
\newcommand{\Succ}{\mathrm{Succ}}
\newcommand{\End}{\mathrm{End}}
\newcommand{\conv}{\operatorname{conv}}
\newcommand{\Rp}{\mathbb{R}^n_+}
\newcommand{\Q}{\mathcal{Q}}
\newcommand{\act}[1]{\textsc{#1}}

\journal{Journal of Multi-Criteria Decision Analysis}

\begin{document}
\begin{frontmatter}
\title{Supplementary Information for:\\ Transition Certification with Composite Indices: The Acceptance Gap Between Compensatory and Noncompensatory Rules}
\author[1]{Amin Abaee}
\ead{amin_abaee@ut.ac.ir}
\affiliation[1]{organization={Independent Researcher}, city={Tehran}, country={Iran}}
\begin{abstract}
This Supplementary Information accompanies the v77 main article. It contains the remaining proofs (S2), the CES diagonal results and exact off-diagonal criterion with finite exploratory checks (S3), the common-shock variants (S4), the normalized-unit biomass illustration (S5), practitioner material (S6), the verification and Lean-coverage record (S7), literature positioning (S8), and a compactness consequence under finite confinement (S9). Main-text results are cited as ``main-text Theorem~2'' etc.; SI results use S-numbering. Exact rational calculations, rational enclosures of selected transcendental comparisons, numerical explorations, and artifacts available for independent reproduction are distinguished in SI~S7 and the main-text Data Availability statement. The Northern-cod value is a historical dimensional comparison only; no parameter is estimated, and no fit, forecast, validation, or empirical performance claim is made. No main-text proof depends on S3--S9.
\end{abstract}
\end{frontmatter}

\section{Abbreviations and the result map}

\begin{table}[!htbp]
\centering\small
\caption{Abbreviations.}
\begin{tabular}{@{}ll@{}}
\toprule
CES & constant elasticity of substitution \\
DFO & Fisheries and Oceans Canada \\
LPI & Living Planet Index \\
LRP & limit reference point \\
MCDA & multi-criteria decision analysis \\
MSY & maximum sustainable yield \\
NCAM & Northern Cod Assessment Model \\
SI & Supplementary Information \\
\bottomrule
\end{tabular}
\end{table}

\begin{table}[!htbp]
\centering\small
\caption{Map between the main-text result numbers of this local v77 edition (numbering retained from the v64 research-article structure) and the technical-report editions v62--v70. Verification and Lean references use earlier numbering; the map is for orientation, with evidence limits stated in SI~S7.}
\begin{tabularx}{\textwidth}{@{}>{\raggedright\arraybackslash}p{0.42\textwidth}>{\raggedright\arraybackslash}p{0.24\textwidth}X@{}}
\toprule
Main text (v77) & Report editions (v62--v70) & Verifier / Lean reference \\
\midrule
Lemma 1 (quantifier inclusion) & Remark 1 & -- \\
Lemma 2 (full cone lossless) & Remark 2 & Lean: full-cone identity \\
Proposition 1 (hierarchy, localisation) & Proposition 3 & grid check 10, 19; Lean \\
Proposition 2 (weight-family monotonicity) & Proposition 4 & -- \\
Lemma 3 (path-to-margin synchronisation) & unnumbered lemma before Theorem 6 (v70) & -- \\
Theorem 1 (finite-menu geometry) & Theorem 6 & -- \\
Remark 1 (fixed-weight countermechanism) & unnumbered remark before Theorem 6 (v70) & grid check 12 \\
Theorem 2 (witnessed separation) & Theorem 5 & grid checks 5--9, 11--22; Lean \\
Proposition 3 (persistence, erasure) & Remark 7, Theorem 8 & grid checks 23--25 \\
Theorem 3 (blend collapse) & Theorem 9 & Lean \\
Proposition 4 (time-sharing) & Proposition 10 & Lean \\
Proposition 5 (rescue threshold) & Proposition 11 & grid check 22 \\
Theorem 4 (substitutability ladder) & Lemma A, Lemma B, Theorems S1, S2 & reported $\sigma$-verifier (57 finite checks); selected Lean rungs (scope in SI~S7) \\
Proposition 6 (common-shock verdict) & Theorems M1, M2, M3 & historically reported 77-check list; artifact limits in SI~S7 \\
S3 Theorem~\ref{thm:offdiag} (off-diagonal criterion) & Theorem G & reported finite scan (50 checks; scope in SI~S7) \\
S4 Theorems S2--S3 (coupled regime, replace-reading collapse) & Theorems M4, CT & historically reported list; see SI~S7 for scope \\
S9 finite-generation consequence under compact confinement & new in v77 & paper proof \\
\bottomrule
\end{tabularx}
\end{table}

\section{Remaining proofs}\label{si:proofs}

\subsection{Persistence under hold-prefix extension (main-text Proposition 3(i))}

\begin{proof}
(Hierarchy.) Backward induction on stages. Base: at the final stage the terminal sets satisfy $G\subseteq G^w\subseteq G^{\mathrm{phys}}$ by main-text Lemma~2 and the definition of the terminal sets. Step: the accepted set of a stage is the set of states from which some action keeps the tube in the safe set and the successor in the next-stage accepted set; applying main-text Proposition~1(i) to the stage's action sets and the induction hypothesis to the next-stage accepted sets preserves the three-way inclusion at the stage level.
(Persistence of strictness.) With \act{hold} the unique prefix action, a prefix state is accepted exactly when it lies in the prefix safe set and in the next-stage accepted set; the stage-0 regions are therefore the witness regions of main-text Theorem~2 pulled back through the (identity) holds, and the strictness witnesses ($I\neq\varnothing$ and the point $(\tfrac12,\tfrac1{10},\tfrac1{10})$) persist.
(Conditionality.) The example below erases the gap because the repair stage's accepted sets cover every first-stage successor while the first-stage tubes remain safe and the original final terminal rules are preserved. A safe continuation cannot repair an earlier tube breach; coincident terminal sets alone do not align path constraints; and coupling actions across stages alone is not sufficient. No general multi-stage erasure rule is asserted.
\end{proof}

\subsection{The erasure datum (main-text Proposition 3(ii))}\label{si:erasure}

\emph{Shared terminal rules.} There are two margins and no reserve. At the end of the final interval the typed terminal set is $G=\{s_1\ge0,s_2\ge0\}$, while the scalarized terminal set at weight $w$ is $G^w=\{w\cdot s\ge0\}$. These are the same rules in the one-stage truncation and the two-stage extension; they are not declared coincident.

\emph{Stage 2, the repair interval.} Its path-safe set is all of $\mathbb R^2$. From every starting state $s$, the single action \act{repair} has tube $\{s\}$ and successor $(1,1)$, which lies in both $G$ and every $G^w$ for $w\in\Wp$. Consequently the stage-2 typed accepted set and every stage-2 scalarized accepted set are all of $\mathbb R^2$.

\emph{Stage 1.} Its path-safe set is the nonnegative orthant, and the two actions have constant tubes at the initial state, with successors
\[
A_1(s)=(s_1-1,s_2+1),\qquad A_2(s)=(s_1+1,s_2-1).
\]
At $z^*=(2/5,2/5)$ both stage-1 tubes are safe. In the two-stage problem, both successors lie in the stage-2 accepted sets $\mathbb R^2$, so each first-stage action is accepted under both doctrines. Backward induction therefore accepts $z^*$ as typed and weak in the extension. No earlier path breach is repaired; the first-stage path is safe and the repair stage enlarges only the continuation accepted sets.

\emph{One-stage truncation with the same terminal rules.} Removing the repair interval leaves $G$ and $G^w$ as the terminal tests after stage 1. Both successors fail the typed terminal test, since their negative coordinate is $-3/5$. For $r=w_2/w_1$ with $w_1,w_2>0$, $A_1$ is scalarized-admissible exactly when $-3/5+(7/5)r\ge0$, i.e. $r\ge3/7$; $A_2$ is admissible exactly when $7/5-(3/5)r\ge0$, i.e. $r\le7/3$. These intervals cover all positive ratios, and the coordinate weights are covered directly by $A_2$ at $r=0$ and $A_1$ at $r=+\infty$. Thus $z^*$ is weak-accepted but typed-rejected in the truncation. The difference is precisely the removal of a repair stage whose accepted sets covered the stage-1 successors, not a change in terminal semantics.

\subsection{The substitutability ladder (main-text Theorem 4)}\label{si:ladderproofs}

Notation as in main-text Section~\ref{main-sec:ladder}: $\lambda_i=1+s_i$, normalised weights $w$, $M_\theta(\lambda;w)=(\sum_iw_i\lambda_i^\theta)^{1/\theta}$, $\sigma=1/(1-\theta)$. The collapse convention is an additional floor-safety rule, not a standard CES identity.

\begin{lemma}[the handshake]\label{lem:handshake}
With weights normalised, $\sum_iw_i\lambda_i\ge1\iff w\cdot s\ge0$; the $\theta=1$ member therefore reproduces the compensatory operator of the main text exactly -- the same acceptance set at every state and weight, and at the witness point the same per-weight plan assignment ($\rho_1=2/3$, $\rho_2=3/2$).
\end{lemma}
\begin{proof}
$\sum_iw_i\lambda_i=\sum_iw_i+\sum_iw_is_i=1+w\cdot s$.
\end{proof}

\begin{lemma}[collapse convention]\label{lem:collapse}
For every $\theta<1$ member, any tube value $\lambda_i\le0$ rejects the plan: a coordinate at or below the collapse level cannot be compensated by surpluses elsewhere.
\end{lemma}
\noindent This rejection rule is imposed uniformly, not derived from every power mean. At $\theta=0$, a zero input with positive weight makes the geometric mean zero; at $\theta<0$, zero is outside the usual power-mean domain. For $0<\theta<1$, however, the mean can remain positive when one coordinate is zero (for example, equal weights on $(0,4)$ give $M_{1/2}=1$), so automatic rejection is an additional safety convention; negative inputs are not standard geometric-mean inputs. The $\theta=1$ member alone is allowed to compensate across collapse: at $(1/2,1,1)$ the fast trough has $\lambda_1=0$ and is rejected at every $\theta<1$ by convention, while the linear family retains its declared compensation rule.

\paragraph{The master equation} On the diagonal $z=(x<1,s,s)$ with $s\in(1,2)$, \act{fast}'s worst-case trough is $\lambda=(s-1,s+1)$ and \act{slow}'s is $(s+1,s-1)$; \act{staged} is inadmissible ($x<1$). Write $p=w_1$ for the weight share, $A=(s-1)^\theta$, $B=(s+1)^\theta$. For $\theta>0$, \act{fast} is accepted at $p$ iff $pA+(1-p)B\ge1$ and \act{slow} iff $pB+(1-p)A\ge1$; both are affine in $p$ and sum to $A+B$. \act{fast}'s acceptance set is a downward interval $[0,p_F]$ (its condition is decreasing in $p$ since $A<B$) and \act{slow}'s an upward interval $[p_S,1]$, with $p_S=1-p_F$ by symmetry; the cover $[0,p_F]\cup[p_S,1]=[0,1]$ holds iff $p_F\ge\tfrac12$, i.e.\ iff the condition holds at $p=\tfrac12$: $(A+B)/2\ge1$. For $\theta<0$ the inequality $M_\theta\ge1$ reads $\sum w_i\lambda_i^\theta\le1$ and the same argument gives $(A+B)/2\le1$. For $\theta=0$ the condition is $(s-1)^p(s+1)^{1-p}\ge1$ with the cover at $p=\tfrac12$: $s^2-1\ge1$. This is the averaging argument: the two plan conditions sum to a weight-independent quantity, so the equal weight is always the binding one.

\begin{proof}[Proof of main-text Theorem 4]
(i) \emph{Nesting.} For fixed positive $\lambda$ and weights, $M_\theta$ is nondecreasing in $\theta$ (the power-mean inequality); the collapse convention preserves this, since acceptance at any $\theta'<1$ forces $\lambda>0$ on the serving plan's tube, where the family is defined. Per plan and weight, hence per weight, hence for the protocol: acceptance at $\theta'$ implies acceptance at every $\theta>\theta'$, so the accepted rungs form a prefix and $\sigma^*(z)$ is well defined. \emph{Finiteness} on gap states with both coordinates in $(1,2)$: at the equal weight, \act{fast}'s condition contains the term $(s_1-1)^\theta$ and \act{slow}'s the term $(s_2-1)^\theta$, both of which tend to $\infty$ as $\theta\to-\infty$ because $s_i-1<1$; so for some finite $\theta$ neither plan serves the equal weight and the state is rejected (on the diagonal this is $(s-1)^\theta\to\infty$). \emph{Positivity}: $s_1+s_2\ge2$ accepts at $\theta=1$. \emph{Only-linear states}: if $\min(s_1,s_2)\le1$, say $s_1\le1$, then \act{fast}'s trough has $\lambda_1=s_1-1\le0$ and is rejected at every $\theta<1$ by Lemma~\ref{lem:collapse}; \act{slow} alone must then cover every weight, including the weight concentrated on coordinate $2$, where its condition is $s_2-1\ge1$, i.e.\ $s_2\ge2$, excluded on a gap state. So only the linear member certifies and $\sigma^*=\infty$.
(ii) Immediate from (i): every gap state is linear-accepted, the accepted rungs form a prefix, and finiteness (or only-linearity) supplies a positive elasticity at which the state is rejected. On the deep-interior subset ($1<s_1<2$, $1<s_2<2$, $s_1+s_2>2$), write $p=w_1\in[0,1]$ and define the linear slacks $L_F(p)=p(s_1-2)+(1-p)s_2$ and $L_S(p)=ps_1+(1-p)(s_2-2)$. The licensing intervals overlap strictly because $s_1+s_2>2$, so $\Delta(p)=\max\{L_F(p),L_S(p)\}>0$ for every $p$. The function $\Delta$ is continuous on the compact simplex $[0,1]$, hence $\eta=\min_{p\in[0,1]}\Delta(p)>0$. All tube inputs $1+s_i-d_{a,i}$ are positive on this deep-interior set. The best-plan CES margin is jointly continuous in $(\theta,p)$ near $\theta=1$ on this compact positive domain; uniform continuity therefore gives an $\varepsilon>0$ such that the cover remains strict for some $\theta\in(1-\varepsilon,1)$. This supplies a finite positive elasticity $\sigma=1/(1-\theta)$ below the linear limit that still accepts the state. On gap states with $\min(s_1,s_2)\le1$, the collapse convention rejects every $\theta<1$ member, so only the linear member accepts.
(iii) Evaluate the master comparison at $s=2$: for $\theta\in(0,1)$ the accepting side is $\ge2$ and $(s-1)^\theta+(s+1)^\theta=1+3^\theta>2$ strictly; for $\theta<0$ the accepting side is $\le2$ and $1+3^\theta<2$ strictly; for $\theta=0$, $4>2$. All strictly accepting at $s=2$, so by continuity some $\varepsilon>0$ has every $s\in(2-\varepsilon,2)$ accepted at that rung, and these are gap states (typed acceptance needs $s\ge2$); the acceptance condition is monotone in $s$ on $(1,2)$ (both $A$ and $B$ move the accepting way), so the accepted set is exactly $[s^*(\theta),2)$ with $s^*(\theta)<2$. The closed forms of the displayed rungs in main-text Table~\ref{main-tab:ladder} follow by direct evaluation ($\theta=2/3$: $s^*=\sqrt{6\sqrt3-9}$; $\theta=1/2$: $\sqrt{s-1}+\sqrt{s+1}\ge2$, equality at $s=5/4$; $\theta=0$: $s^2\ge2$; $\theta=-1$: $(s^2-1)/s\ge1$; $\theta=-2$: $(s+1)^2+(s-1)^2\le2(s^2-1)^2\iff s^2\ge3$; $\theta=-m$: $(s-1)^{-m}+(s+1)^{-m}\le2$). Nesting gives nondecreasing critical floors as elasticity falls; the strict inequalities among the displayed rungs follow from these explicit roots and brackets, not from nesting alone. At $\theta\to-\infty$, the power mean tends to the minimum over the positive-weight support. A strictly positive full-support weight therefore requires every coordinate to meet its floor along the tube, while typed safety passes every support minimum; the Leontief-limit all-weight accepted set is exactly the typed set. Hence $\inf_z\sigma^*(z)=0$ and only $\sigma=0$ is uniformly safe.
\end{proof}

\subsection{The common-shock verdict (main-text Proposition 6)}\label{si:csproofs}

The variant is defined in main-text Section~\ref{main-sec:commonshock} and SI~S4 below. Write $A(p)=ps_1+(1-p)s_2$ for the aggregate margin at weight share $p$.

\begin{proof}[Proof sketch]
(i) \act{fast}'s common-class requirement is $A(p)\ge\max(2p,\ p+\tfrac12)$ (its own tent $3/2$ in $s_1$ plus the $s_1$-event, and the $s_2$-event weighted by $1-p$); \act{slow}'s is the mirror; \act{staged}'s ($x\ge1$) is $A(p)\ge\max(p,1-p)/2-1/8$, the event exposure of a growing floor (kinked in $p$; its admissibility decomposes exactly into the two affine conditions of the kink's branches). The pointwise minimum of the two principal requirements is a tent with its only kink at $p=\tfrac12$, so an affine $A$ covers it exactly at the three abscissae $p\in\{0,\tfrac12,1\}$; the closed forms follow: typed acceptance needs $s_1\ge2$ with $s_2\ge\tfrac12$ (or the mirror) or the staged branch $x\ge1$, $\min(s)\ge\tfrac38$; weak acceptance with $x<1$ needs $s_1\ge\tfrac12$, $s_2\ge\tfrac12$ (the endpoints $p=0,1$) and $s_1+s_2\ge2$ (the kink). The licensing endpoint of \act{fast} on the surviving region equals its action-indexed endpoint $1/(1+\rho_1)$, since the binding branch there is the own-tent branch. The closed forms above follow from the stated pointwise maxima and the three critical abscissae. A prior report lists finite grid checks, but the corresponding checker and run log were not found in the inspected versioned files; no v77 claim relies on that unavailable artifact.
(ii) The augmented menu adds $\act{staged}_\kappa$, admissible at $x+\kappa\ge1$ with the same floor requirement $s_i\ge\tfrac38$; $\kappa$ cannot move the floors, so on $\{\min(s)<\tfrac38\}\setminus\Vtyp^{\mathrm{cs}}$ no finite $\kappa$ certifies. On the complement the requirement is the action-indexed one, $\kappa=(1-x)_+$. The $3/8$ boundary is structural: the event exposure of a growing floor is $\max(0,\delta-1/8)=3/8$ at $\delta=1/2$.
(iii) The blend's $s_1$-trough is $s_1-\delta-\delta_{\mathrm{blend}}\cdot\tfrac32$ and its $s_2$-trough $s_2-\delta-(1-\delta_{\mathrm{blend}})\cdot\tfrac32$ (each coordinate takes its event plus the corresponding fraction of the principal tent); the window is the nonnegativity interval of the two intersected with $[0,1]$, nonempty iff $s_1\ge\delta$, $s_2\ge\delta$ and $s_1+s_2\ge\tfrac32+2\delta$.
\end{proof}
